\documentclass[../../../../main.tex]{subfiles}
\begin{document}

\begin{flushright}
\footnotesize\textit{【自動文字起こし・要確認】}
\end{flushright}

[解] $n$ 個から $2k$ 個をとり出し, これを $2$ つずつの組に分ける場合の数と同じである。

したがって, $n \ge 2k$ のもとで
\[
\frac{{}_n\mathrm{C}_{2k} \cdot {}_{2k}\mathrm{C}_2 \cdot {}_{2k-2}\mathrm{C}_2 \cdots {}_4\mathrm{C}_2 \cdot {}_2\mathrm{C}_2}{k!}
\]
\[
= \frac{\frac{n!}{2k!(n-2k)!} \cdot \frac{2k(2k-1)}{2} \frac{4\cdot 3}{2} \cdots \frac{2\cdot 1}{2}}{k!}
\]
\[
= \frac{n!}{k!(n-2k)!} \cdot \frac{1}{2^k} \quad \text{通り}
\]
であり, $n < 2k$ の時は $0$ 通り。以上まとめて,
\[
f(n, k) = 
\begin{cases}
\frac{n!}{(n-2k)! \cdot k!} \frac{1}{2^k} \, \text{通り} & (n \ge 2k) \\
0 \, \text{通り} & (n < 2k)
\end{cases}
\]

次に, $f(n, k) = f(n, 1)$ について, $n \ge 2$ から $f(n, 1) \neq 0$ だから, $n \ge 2k$ のもとで考えれば良い。
\[
f(n, k) = f(n, 1)
\]
\[
\iff \frac{1}{2^k} \frac{n!}{k!(n-2k)!} = \frac{1}{2} \frac{n!}{(n-2)!}
\]
\[
\iff 2^{k-1} k! (n-2k)! = (n-2)!
\]
\[
\iff 2^{k-1} = \frac{(n-2)(n-3) \dots (n+1-2k)}{k!} \quad \dots \text{①}
\]
ここで, 連続した整数は $k!$ で割り切れるから, $A \in \mathbb{N}$ として,
\[
(n-2) \cdots (n+1-2k) = A(2k-2)!
\]
とおける。①から
\[
\text{①} \iff 2^{k-1} = A \cdot (2k-2)(2k-3) \cdots (k+1) \quad \dots \text{②}
\]
同様に, $(2k-2) \dots (k+1)$ は $(k-2)!$ で割り切れるから, $k \ge 5$ の時 $2^{k-1}$ が $3$ で割り切れることになり不適。又, $k=2, 4$ の時は ②から, $2^{k-1}$ が各々 $3$ ($k=2$), $5$ ($k=4$) で割り切れることになり不適。以上から $k=3$ が必要。この時①から,
\[
3 \cdot 2 \cdot 2^2 = (n-2) \cdots (n-5)
\]
\[
1 \cdot 2 \cdot 3 \cdot 4 = (n-2) \cdots (n-5)
\]
となり, $n=6$ とすれば十分。以上から, もとめるのは $(n, k) = (6, 3)$ である。

\end{document}
